Al taller de l'institut dissenyem un pèndol amb un suport i un pes que penja d'un fil de longitud $L$ i de massa negligible. Quan prenem mesures del període $T$ d'aquest pèndol en funció de la longitud del fil, obtenim les dades següents:

\begin{center}
\renewcommand{\arraystretch}{1.4}
\begin{tabular}{|l|*{5}{p{1.3cm}|}}
\hline
\textit{Longitud $L$ (m)} & 0,20 & 0,30 & 0,50 & 0,80 & 1,20 \\
\hline
\textit{Període $T$ (s)} & 0,90 & 1,08 & 1,43 & 1,75 & 2,20 \\
\hline
\textit{$T^2$ (s$^2$)} & & & & & \\
\hline
\end{tabular}
\end{center}

\begin{apartat}{1,25}
Ompliu la tercera fila de la taula anterior amb els valors del període al quadrat. Representeu $T^2(L)$ a la quadrícula de sota i justifiqueu per què els punts estan alineats. Calculeu el pendent d'aquest gràfic. A partir del resultat, calculeu també la intensitat del camp gravitatori.
\end{apartat}

\figura[0.6]{fig1.png}

\begin{apartat}{1,25}
Deixem caure el pèndol d'1,50~kg de massa des d'un angle de $30^\circ$ respecte a la vertical i considerem negligible el fregament amb l'aire. Tenint en compte que la longitud del pèndol és d'1,00~m, calculeu les energies potencial i cinètica màximes i l'energia mecànica del pèndol, així com la velocitat màxima que assolirà. Trieu com a referència d'energia potencial zero la posició d'equilibri.
\end{apartat}

\dades{El pèndol descriu un moviment harmònic simple de període $T = 2\pi\sqrt{\dfrac{L}{g}}$.}
